\chapter{第四版：第0章§§16--17的译校与核查}
\section{范围、证据与分类}
本轮连续译文的底本为用户上传的 \texttt{EGAIV.1.pdf}，第0章§16“诺特局部环中的维数与深度”和§17“正则环”，从（16.1.1）至（17.3.9）及其后的两段文字。它们位于该文件第21--50页的相应部分，原分册印刷页码为22--51。还核对了 \texttt{EGAIV.4.pdf} 第344、346页的相关作者勘误（印刷第345、347页）。原扫描同时出现分册页码与期刊连续页码，本报告统一采用分册印刷页码，并辅以PDF页序。

本轮编号分为：D4-A（作者正式勘误）、D4-E（有明确反例的原字面断言）、D4-G（可修补的证明步骤）、D4-P（仍需补足的证明事项）、D4-T（文字、对象或引文调整）、D4-H（假设、定义及阅读边界）。\textbf{编号总数不是错误定理的数量。}尤其是D4-E02只是零模边界，D4-A02是作者删去多余假设的推广，D4-H中多项只是防止误读。

译文按初版组织与详略排列。显式采用作者改正之处有脚注；仍保留初版原句之处，也在紧邻脚注中说明修订。下面的反例与补证是译校者的分析，不伪装成原著证明；未逐行复核此前合订本全部内容，历史报告也未在本轮获得再次独立认证。

\section{作者正式勘误：五项}
\audit{D4-A01}{（16.1.4）的悬链性表述缺少整环限定}{作者已正式更正，不是本次首次发现}
\source{IV第一部分印刷第23页；IV第四部分印刷第347页，Err$_{\mathrm{IV}}$，5}
初版将悬链性与所有 $\qq\subset\pp$ 都满足
\[
 \dim A_{\qq}+\dim(A_{\pp}/\qq A_{\pp})=\dim A_{\pp}
\]
直接等同。作者要求把“说 $A$ 是悬链的”改为“说\textbf{一个整环} $A$ 是悬链的”。

区别在于：悬链性比较\emph{同一对端点之间}的饱和链，不保证不同极小素理想到同一极大素理想的链等长。可以用
\[
 A=\bigl(k[x,y,z]/(xy,xz)\bigr)_{(x,y,z)}
\]
作具体检验。令 $\mathfrak p=(x)$、$\mathfrak q=(y,z)$、$\mathfrak m=(x,y,z)$ 表示商环中的理想。两个不可约分支分别为
\[
 A/\mathfrak p\simeq k[y,z]_{(y,z)},\qquad
 A/\mathfrak q\simeq k[x]_{(x)},
\]
维数为2和1，且只在闭点相交。每个固定端点间的素理想区间位于相应分支；前一分支中极小到极大的饱和链长度为2，后一分支中为1，各自满足悬链性，故 $A$ 悬链。但
\[
 \dim A_{\mathfrak q}=0,\qquad
 \dim(A_{\mathfrak m}/\mathfrak qA_{\mathfrak m})=1,
 \qquad\dim A_{\mathfrak m}=2.
\]
所写等式变成 $0+1=2$。这与作者增加整环条件一致，不应继续引用未限定的初版句子。

\audit{D4-A02}{（16.1.9）的诺特假设可以删去}{作者推广；初版较强假设下的结论不因此错误}
\source{IV第一部分印刷第25页；IV第四部分印刷第347页，Err$_{\mathrm{IV}}$，6}
初版用 $A$ 诺特保证 $\End_A(M)$ 及其子模 $B_M$ 有限型。作者指出可不用这一途径：任意 $u\in B_M$ 在有限型 $A$-模 $M$ 上作用，行列式技巧给出 $u$ 的整性。具体地，若 $m_1,\ldots,m_r$ 生成 $M$，写
\[
 um_i=\sum_j a_{ij}m_j.
\]
伴随矩阵恒等式表明，首一多项式 $\det(TI-(a_{ij}))$ 在 $T=u$ 处零化 $M$。由于 $B_M$ 就是 $B$ 在 $\End_A(M)$ 中的像，作用忠实，故该多项式在 $B_M$ 中确实为零。所以 $B_M$ 整于 $A_M$，再应用（16.1.5）。

这不是声称任意环上有限型模的端同态环必有限型，也不是将初版“子模有限型”的论证原样移到非诺特情形。

\audit{D4-A03}{（16.2.2.2）交滤过的母模误植}{作者已正式更正的公式错误}
\source{IV第一部分印刷第26页；IV第四部分印刷第345页，勘误A部}
初版印为 $M'_n=M'\cap\mathfrak q^nM'$；应为
\[
 \boxed{M'_n=M'\cap\mathfrak q^nM.}
\]
此时才有正合列
\[
 0\longrightarrow M'/(M'\cap\mathfrak q^nM)
 \longrightarrow M/\mathfrak q^nM
 \longrightarrow M''/\mathfrak q^nM''\longrightarrow0,
\]
并能将长度等式与Artin--Rees良好滤过结合。取 $A=k[[t]]$、$M=A$、$M'=tA$、$\mathfrak q=(t)$，则 $M'\cap t^nM=t^nA$（$n\geq1$），而 $t^nM'=t^{n+1}A$，两者并不相同。正文采用作者修订，并注明初版原式。

\audit{D4-A04}{（16.5.12）不能直接套用未限定的（16.1.4.2）}{作者已给出证明修补}
\source{IV第一部分印刷第39页；IV第四部分印刷第347页，Err$_{\mathrm{IV}}$，7}
作者指出，关系（16.5.11.1）等价于将 $A$ 换成各极小素理想商环 $A/\mathfrak r$ 后的相同关系。因此可先归约到整环，再使用更正后的（16.1.4.2）。译文保留初版证明，并以脚注指向该修补及书末选译。这里\textbf{没有否定Cohen--Macaulay局部环及有关商环的悬链性结论}。

\audit{D4-A05}{（17.3.7）误引不存在的（17.1.10）}{作者已正式更正的引文}
\source{IV第一部分印刷第50页；IV第四部分印刷第345页，勘误A部}
多项式环正则性的证明应使用（17.3.3），不是初版印出的（17.1.10）。相应步骤是：$A_{\mathfrak p}\to B_{\mathfrak q}$ 为平坦局部同态，底环正则，闭纤维局部环正则，故总环正则。本版正文按作者勘误更换引文。

\section{有明确反例的两项字面断言}
\audit{D4-E01}{（16.4.2）对任意模的极大正则序列判据不成立}{缺少有限型或相应有限伴随素理想假设}
\source{IV第一部分印刷第32--33页，上传PDF第31--32页}
初版只假定 $A$ 为诺特局部环，$M$ 为任意 $A$-模，却声称：由极大理想中的元素组成的 $M$-正则序列 $(x_1,\ldots,x_n)$ 极大，当且仅当
\[
 \Hom_A\bigl(k,M/(x_1M+\cdots+x_nM)\bigr)\ne0.
\]
其证明用“极大理想包含于伴随素理想之并，故等于其中之一”。对无限族，不能这样使用有限素理想回避。

\subsection*{一个完全明确的反例}
取任意域 $k$，令
\[
 A=k[x,y]_{(x,y)},\qquad\mathfrak m=(x,y)A,
 \qquad
 M=\bigoplus_{\substack{\mathfrak p\in\Spec A\\\htp\mathfrak p=1}}A/\mathfrak p.
\]
这是诺特局部环上的一个无限直和模。下面分别核对两边。

任取 $0\ne f\in\mathfrak m$。这个环是多项式唯一分解环的局部化；$f$ 有一个非单位不可约因子 $g$。令 $\mathfrak p=(g)$，它是高度1素理想，并且 $f\in\mathfrak p$。于是，在直和因子 $A/\mathfrak p$ 中，非零元 $1+\mathfrak p$ 被 $f$ 零化。$f=0$ 的情形更直接。因此，$\mathfrak m$ 中不存在任何 $M$-正则元素，空正则序列不能继续延长，因而极大。

另一方面，每个 $A/\mathfrak p$ 都是一维局部整环。由于 $\mathfrak p\ne\mathfrak m$，可取 $u\in\mathfrak m-\mathfrak p$；其在 $A/\mathfrak p$ 中非零，数乘单射。所以
\[
 (0:_{A/\mathfrak p}\mathfrak m)=0.
\]
直和中的元素只有有限个非零分量；一个元素若被 $\mathfrak m$ 的全部元素零化，则其每个分量均被 $\mathfrak m$ 零化，因而全为零。故
\[
 \Hom_A(k,M)=(0:_M\mathfrak m)=0.
\]
这就得到“正则序列极大，但所写Hom为零”，直接否定任意模版本的必要性。

\subsection*{有限型情形的修补与影响范围}
若 $M$ 有限型，令 $N=M/(x_1M+\cdots+x_nM)$；则 $N$ 有限型，$\Ass N$ 有限。序列极大意味着 $\mathfrak m$ 的每个元素在 $N$ 上都是零因子，故
\[
 \mathfrak m\subset\bigcup_{\mathfrak p\in\Ass N}\mathfrak p.
\]
有限素理想回避给出 $\mathfrak m\subset\mathfrak p$ 对某个 $\mathfrak p\in\Ass N$ 成立。由于 $\mathfrak m$ 极大，$\mathfrak p=\mathfrak m$，所以 $\Hom_A(k,N)\ne0$。反方向对任意模都成立：若 $0\ne v\in N$ 被 $\mathfrak m$ 零化，则任何 $x\in\mathfrak m$ 都不能在 $N$ 上单射。

因此，（16.4.4）等显式要求有限型的深度结论不被上述反例推翻。（16.4.3）的证明只使用“非极大序列的最后商模没有非零 $\mathfrak m$-零化元”这一正确方向，也不能据此宣布（16.4.3）错误。

\audit{D4-E02}{（17.2.7）的射影维数递增式漏排零模}{边界条件遗漏；非零有限型模版本成立}
\source{IV第一部分印刷第44页；射影维数约定见印刷第42页（17.2.2（i））}
初版在（17.2.6）的假设下写
\[
 \pd_A(M/xM)=\pd_A M+1
\]
却没有注明 $M\ne0$。本书的约定使零模为射影模，故 $\pd_A0=0$；本书的正则元素定义只要求数乘单射，所以每个 $x$ 对零模都正则。取 $A=k[[t]]$、$x=t$、$M=0$，左端为0，右端为1。

修补是加上 $M\ne0$。若 $M$ 非零有限型，Nakayama引理给出 $M/\mathfrak mM\ne0$。原文得到的正合列
\[
 0\to\Tor_i^A(M,k)\to\Tor_i^A(M/xM,k)
 \to\Tor_{i-1}^A(M,k)\to0
\]
此时能正确检测最高非零Tor。若 $\pd_A M=d<\infty$，则 $\Tor_d^A(M,k)\ne0$，从而商模的射影维数为 $d+1$。若 $\pd_A M=\infty$，对任意大的次数仍存在非零Tor，以上正合列使商模也有无限射影维数。故非零版本并未被反驳。

\section{证明步骤与待补事项}
\audit{D4-G01}{（16.2.3.2）归纳中多写了尚未得到的等号}{可用所需不等式完成原证明}
\source{IV第一部分印刷第27页}
这一阶段正在证明 $\dim M\leq d(M)$，尚未完成三者相等的定理。原文却以归纳假设写
$d(N/xN)=\dim(N/xN)$。实际只需
\[
 n-1\leq\dim(N/xN)\leq d(N/xN)\leq d(N)-1.
\]
第二个不等式正是较小Hilbert--Samuel次数下的归纳假设，第三个来自 $x$ 在 $N$ 上正则及（16.2.3.1（iii））。所以 $n\leq d(N)$，不必预先使用完整等式。这是局部的证明表述问题，不是否定Krull--Chevalley--Samuel定理。

\audit{D4-P01}{（16.3.8）半局部环中所选元素未保证在根基内}{证明的选取步骤需补；本轮未给出原命题反例}
\source{IV第一部分印刷第31页}
原命题要选出参数系的一部分，依（16.3.6）这些元素必须在Jacobson根基 $\mathfrak r$ 中。证明先选 $x\in\mathfrak a$ 避开所有极小素理想，却未保证 $x\in\mathfrak r$，随后使用含根基条件的（16.3.4）。

这一限制不是自动满足的。令 $R=k[[t]]$、$S=k[[u,v]]$，取
\[
 A=R\times S,\quad\mathfrak a=(t)R\times S,\quad x=(t,1).
\]
$A$ 是诺特半局部环，$\mathfrak a\ne A$，且 $\operatorname{ht}\mathfrak a=1$。它的两个极小素理想是 $0\times S$ 与 $R\times0$；$x$ 避开两者，却不在
$\mathfrak r=(t)R\times(u,v)S$ 中。并且
\[
 \dim A=2,\qquad A/xA\simeq k,\qquad\dim(A/xA)=0\ne\dim A-1.
\]
所以证明按字面允许的这种选择，既不能成为参数系的一部分，也不满足所用的降维等式。

这\textbf{不反驳存在性命题}：在同一例子中改取 $(t,u)$ 即能满足所需性质。一般半局部情形还需同时控制各极大点下的选择与余维条件；本轮保留原命题，明确标记该步待补，未将一个仅在局部情形直接成立的证明冒充一般修补。

\section{文字、对象与引文调整}
\audit{D4-T01}{（16.1.1）维数定义的引用}{原号与定义位置不符}
初版引用（14.2.1），该条定义余维；组合维数的定义是（14.1.2）。正文保留原引，脚注指出。位置：印刷第22页。

\audit{D4-T02}{引理编号（16.12.2.1）}{初版标题误排，后文均按16.2.2.1引用}
初版印刷第26页标题多出数字1。本译稿排作（16.2.2.1），记录原号，不补造一条所谓“16.12.2.1”。

\audit{D4-T03}{（17.2.5）Ext局部化应解消第一变量}{证明中的模名错置}
\source{IV第一部分印刷第43页}
原文要证明
\[
 S^{-1}\Ext_A^i(N,M)\simeq
 \Ext_{S^{-1}A}^i(S^{-1}N,S^{-1}M),
\]
其中 $A$ 诺特、$N$ 有限型，$M$ 可任意。这里应取 $N$ 的有限型自由左解消，而不是原文所说的 $M$ 的自由解消。若 $P_\bullet\to N$ 为这样的解消，则逐项有
\[
 S^{-1}\Hom_A(P_j,M)
 \simeq\Hom_{S^{-1}A}(S^{-1}P_j,S^{-1}M).
\]
局部化正合，故与取同调交换，得到所需同构。初版的结论保留，译注指出对象错置。

\audit{D4-T04}{（17.2.13）的自由解消方向}{Tor计算应为左解消}
初版印刷第45页两次写“右解消”。用于计算 $\Tor_i^A(M,N)$ 的自由解消应为
$\cdots\to L_1\to L_0\to M\to0$。正文明确采用左解消并说明原字样；这是方向文字错误，不是平坦基变换结论被推翻。

\audit{D4-T05}{（17.3.4）深度降1公式的引用}{应引16.4.6（ii），不是（i）}
初版印刷第49页在取 $M$-正则元素后写 $\prof(M/xM)=\prof M-1$，却引（16.4.6（i））。正确位置为（ii），原证明路线不受影响。

\audit{D4-T06}{（17.3.9）被表示为商环的对象漏了局部化}{应为 $S^{-1}C$，不是 $C$}
\source{IV第一部分印刷第51页}
若 $C=B'/I$，$S'$ 是 $S$ 在 $B'$ 中的逆像，正确关系是
\[
 S^{-1}C\simeq S'^{-1}B'/S'^{-1}I.
\]
初版相应句子的主语写成 $C$。取 $B'=C=\mathbf Z$、$I=0$、$S=\mathbf Z-\{0\}$，$S'^{-1}B'=\mathbf Q$。不存在从 $\mathbf Q$ 到 $\mathbf Z$ 的保单位环同态，因为2的逆元将不得不映到 $\mathbf Z$ 中2的逆元。所以 $\mathbf Z$ 不是 $\mathbf Q$ 的商环；而 $S^{-1}C=\mathbf Q$ 正是所需商环。正文采用正确对象并保留原式说明。

\section{假设、定义与阅读边界}
\audit{D4-H01}{（16.1.4）的“双等维”}{继承早期译注，不把有问题的等价条件重新当作定义}
应结合（14.3.3）的核查阅读，尤其不能把等维、等余维、悬链三个条件的合取，无条件等同于全部极大链等长。本轮没有重新验证历史报告的每个反例；相关提醒只防止把早期已标注的缺口带入新节。

\audit{D4-H02}{（16.2.1）定义理想的半局部条件}{“商环Artin”等价说法须保留根基包含条件}
原文先要求 $\mathfrak q\subset\mathfrak r$ 且包含 $\mathfrak r$ 的某个幂。若只保留“$A/\mathfrak q$ 是Artin环”，不够：取
$A=k\times k$、$\mathfrak r=0$、$\mathfrak q=k\times0$，则商环是域，然而 $\mathfrak q\not\subset\mathfrak r$。因此此处应理解为在前述根基包含限制下的等价刻画。

\audit{D4-H03}{（16.2.3.1（iii））左下标不能漏}{是 $\Ker x$ 的Hilbert--Samuel多项式，不是 $M$ 的}
原式为
\[
 P_{\mathfrak q}({}_xM,n)-P_{\mathfrak q}(M/xM,n),
 \qquad{}_xM=\Ker(x:M\to M).
\]
以 $A=k[[t]]$、$M=A$、$x=t$、$\mathfrak q=(t)$ 检验：$d(M)=1$，${}_xM=0$，正确差为 $-1$，次数0；若漏掉左下标，错误差变成 $n-1$，次数1，不满足至多 $d(M)-1$。本次按扫描页排出左下标，不把文本识别遗漏当成原著错误。

\audit{D4-H04}{（16.3.6）参数系与定义理想的作用环}{应区分在 $A$ 中与在 $A/\Ann M$ 中}
取 $R=k[[t]]$、$A=R\times k$、$M=R$（经第一投影作用），则
\[
 \mathfrak r=(t)R\times0,\quad
 \Ann_A M=0\times k,\quad\dim M=1.
\]
元素 $x=(t,0)$ 构成 $M$ 的参数系，因为 $M/xM\simeq k$。但
\[
 Ax+\Ann M=(t)R\times k\not\subset\mathfrak r,
\]
不是原文先前定义的 $A$ 的定义理想。在 $A/\Ann M\simeq R$ 中，其像则是定义理想。也可像（16.2.3.3）那样改用
$Ax+(\mathfrak r\cap\Ann M)$。这不否定“参数系可在零化子商环中理解”的结论，只须写清作用环。

\audit{D4-H05}{（16.4.1）的零模边界}{涉及参数系长度时需非零模语境}
原文随后明确以 $M\ne0$ 得到正则序列长度至多 $\dim M$；按本书 $\dim0=-\infty$ 而零模上所有数乘均单射，故前面的“参数系”说法不能机械套到零模。译文保留原范围，并将此列作边界提醒。

\audit{D4-H06}{（16.4.3）不受D4-E01反例直接否定}{只用了前一引理的正确方向}
若 $\Hom_A(k,N)\ne0$，则存在被 $\mathfrak m$ 零化的非零元，所以正则序列不能继续延长。取逆否命题，“仍能延长”推出 $\Hom_A(k,N)=0$。这个方向不要求有限型，足以支持原文该步归纳。不能因为引理某个方向失效就把所有引用它的后续命题都列为错误。

\audit{D4-H07}{正则参数系与正则序列构成的参数系}{两者不应混译}
在 $k[[t]]$ 中，$(t^2)$ 是参数系，也是正则序列，但不生成极大理想 $(t)$，所以不是（17.1.6）的正则参数系。这一区分在原文已经强调，不是本次发现的新错误。

\audit{D4-H08}{（17.2.11）Tor判据的有限型环节}{先对有限型模应用17.2.6，再用17.2.8}
由剩余域射影维数至多 $n$ 可知任意模 $M$ 的高次 $\Tor_i(M,k)$ 消失；但（17.2.6）反推射影维数时要求 $M$ 有限型。正确读法是先取有限型 $M$，推出它们的射影维数均至多 $n$，再由（17.2.8）的条件a'）推出所有模的全局维数界。结论无须修改。

\audit{D4-H09}{（17.2.15）的平坦态射}{应用所引命题需要局部茎同态}
（17.2.13）对平坦\emph{局部}环同态作出结论。因而（17.2.15）应在局部赋环空间及局部态射的约定下理解，不能将“茎环都是局部环”当成“茎同态为局部同态”。本轮未因作者省略重申约定而判定该命题错误。

\audit{D4-H10}{（17.3.1）归纳量与全局维数重用字母}{显式更换记号，不改变原证明}
原证明先以 $n=\rk_k(\mathfrak m/\mathfrak m^2)$ 归纳，随后第一种情形又写 $n=\gldim A$。本译稿将后者改记 $d$，并注明。它有限是因为 $\pd_A\mathfrak m<\infty$、$0\to\mathfrak m\to A\to k\to0$ 及（17.2.11），不是预先假定要证明的正则性。

\audit{D4-H11}{§17末尾关于“重要的环”的叙述}{保留历史语境，不改写成无条件定理}
原著最后一句关于代数几何应用中“一切重要的环”都是正则环商环，是叙述性概括；它没有给出“重要”的数学定义。本版保留原句，但不将其扩张成任意环的普遍结论。

\section{外部交叉核验的限定}
本轮译文以用户上传法文原著为底本。另查看Stacks Project的§10.72（标签00LE）、§10.68（标签0AUH）及§10.110（标签065U），仅用来对照有限型假设、零模约定和正则局部环的全局维数结论。其深度论的关键Ext刻画明确要求非零有限型模；其正则序列定义还要求最终商模非零，与本书初版约定不同。因此不能把现代定义静默替换进EGA译文，也不能用定义差异消除D4-E02在本书自身约定下的边界问题。上述网页不是这两节的翻译来源。

\section{本轮审校结论与后续接续点}
本轮已形成第0章§§16--17的连续译稿，保留初版各编号条目、证明、公式和原注，并补充相关作者勘误与译校注。编译、条目清点、单文件独立编译及页面检查的实际结果，以工程包 \texttt{checks/} 中本次记录为准；不以可编译性替代数学正确性审定。

下一连续译文应从IV第一部分印刷第51页的§18章首及（18.1.1）开始。第0章§§18--23、第IV章主体，以及第五卷预备笔记余文尚未完成。本轮没有把这些材料的原文页、目录或提要计入已译正文。
